\input{../tex-inputs/latex-leseansicht-vorspann}

\section[Arthur Schnitzler an Gustav Schwarzkopf, 18. 9. 1899]{L04423 Arthur Schnitzler an Gustav Schwarzkopf, 18. 9. 1899}
\nopagebreak\mylabel{L04423v}
\rehead{ }\normalsize\beginnumbering\briefempfaengerindex{Schwarzkopf, Gustav@\textsc{Schwarzkopf, Gustav}!zzzSchnitzler, Arthur@\emph{von Arthur Schnitzler}!1899-09-182@{18. 9. 1899}|(be}
\toendnotes[C]{\smallbreak\pagebreak[2]}
\correspDesc{Versand  durch Arthur Schnitzler am 18. 9. 1899 in Nürnberg
\newline{}Übermittlung  am 18. 9. 1899 in Nürnberg
\newline{}Zustellung  am 20. 9. 1899 in Wien
\newline{}Erhalt  durch Gustav Schwarzkopf am 20. 9. 1899 in Wien}\toendnotes[C]{\smallbreak}
\Standort{DLA, A:Schnitzler, HS.1985.1.1897.}
\physDesc{Bildpostkarte, 71 Zeichen
\newline{}Handschrift: Bleistift, deutsche Kurrent
\newline{}Versand: 1) Stempel: »\nobreak{}\oindex{Nürnberg@\textbf{Nürnberg}|pwk}Nuernberg, 18 Sep \textcolor{gray}{99}, 3–4NM\nobreak{}«.   2) Stempel: »\nobreak{}\oindex{I., Innere Stadt@\textbf{I., Innere Stadt}|pwk}Wien 1/1 1, 20 9 \textcolor{gray}{99}, 8–9½V., Bestellt\nobreak{}«. }\pstart{}{\pb}Hrn \textsc{Gustav Schwarzkopf}\pend{}\pstart{}Wien\oindex{Wien@\textbf{Wien}|pw}\pend{}\pstart{}\textsc{I. Tiefer Graben 23}\oindex{Wien@\textbf{Wien}!I., Innere Stadt@\textbf{I., Innere Stadt}!Tiefer Graben 23@\textbf{Tiefer Graben 23}|pw}.\pend{}{\bigskip}
\pstart
           \centering{}{\pb}\textcolor{gray}{\textbf{Nürnberg, Burg\oindex{Nürnberger Burg@\textbf{Nürnberger Burg}|pw} vom Spittlerthor\oindex{Spittlertor@\textbf{Spittlertor}|pw}.}}\pend
           \vspace{1em}
\pstart
           \noindent{}{\pb}Herzliche Grüße!\pend
           \pstart Ihr \spacefill\mbox{Arthur}\pend{}\selectlanguage{ngerman}\endnumbering\briefempfaengerindex{Schwarzkopf, Gustav@\textsc{Schwarzkopf, Gustav}!zzzSchnitzler, Arthur@\emph{von Arthur Schnitzler}!1899-09-182@{18. 9. 1899}|)be}\mylabel{L04423h}
\begin{anhang}
\end{anhang}\newcommand{\dateiname}{L04423}\newcommand{\titel}{Arthur Schnitzler an Gustav Schwarzkopf, 18. 9. 1899}\newcommand{\editorInnen}{Herausgegeben von Jahnke, SelmaMüller, Martin Anton}\input{../tex-inputs/latex-leseansicht-abspann}
